\subsection*{K 敏感性与外生风险补充检验}

正文“模型建立与求解”章只保留参数排序、尾部风险和外生风险摘要。附录在短期模型检验之后补充长期敏感性和外生风险证据，主要用于说明扰动尺度从何而来、外部风险变量如何进入解释，以及 2026 年冲突窗口相对历史窗口是否具有特殊性。

\appendixfigunit{0.38\textheight}{历史波动参数校准}{长期蒙特卡洛扰动不应凭经验设置，因此先用 2017--2025 年同长度历史窗口估计波动、最大回撤和误差分布。该图给出扰动尺度的可追溯来源。}{0.96}{0.28\textheight}{历史波动参数校准.png}{附图 7 历史波动参数校准}

附图 7 的作用是把长期随机扰动落回历史经验区间：模型并非任意扩大波动，而是用同长度窗口的收益波动和回撤分布约束后续情景树的扰动尺度。

\appendixfigunit{0.49\textheight}{传统蒙特卡洛路径云}{该图展示不引入状态转移结构时，多参数联合扰动形成的价格路径云。它作为长期情景树的朴素随机模拟参照，用来区分普通参数扰动和状态切换风险。}{0.98}{0.32\textheight}{传统蒙特卡洛路径云图.png}{附图 8 传统蒙特卡洛路径云图}

\appendixfigunit{0.40\textheight}{GPR 地缘风险指数滞后审计}{GPR 不直接替代附件价格，也不用于同日解释同日价格。这里检验滞后风险指数与布伦特收益的关系，确认其更适合作为风险环境变量。}{0.88}{0.23\textheight}{地缘风险指数滞后审计.png}{附图 9 地缘风险指数滞后检验}

\appendixfigunit{0.36\textheight}{滞后 GPR 与月度收益}{该散点关系说明，滞后 1 月 GPR 与布伦特月度收益的线性解释力有限。因此正文只将 GPR 用作长期风险溢价和状态转移概率的约束变量。}{0.78}{0.21\textheight}{滞后GPR与布伦特收益关系.png}{附图 10 滞后 1 月 GPR 与布伦特月度收益关系}

\appendixfigunit{0.40\textheight}{OVX 隐含波动率滞后检验}{OVX 反映市场对原油期权波动率的定价预期，适合刻画不确定性和尾部风险。该图用于确认 OVX 主要约束长期波动和风险状态，而不是直接给出油价涨跌方向。}{0.96}{0.29\textheight}{OVX隐含波动率滞后检验.png}{附图 11 OVX 隐含波动率滞后检验}

\appendixfigunit{0.40\textheight}{历史同长度窗口极端性}{将 2026 年冲突窗口放回 2017--2025 年同长度窗口中比较，可以判断本题样本是否处于异常高涨幅、高波动或高预测难度区间。}{0.96}{0.31\textheight}{历史窗口极端性检验.png}{附图 12 冲突窗口历史极端性}

附图 11 和附图 12 分别从市场定价风险与历史窗口位置说明同一件事：2026 年冲突窗口不是普通平稳波动段，因此长期模型需要显式保留状态切换和尾部区间，而不能只给一条确定性均值线。

\appendixfigunit{0.42\textheight}{历史基准误差分布}{该图比较历史同长度窗口上的基准误差分布，用于说明 2026 年冲突窗口相对普通市场阶段的预测难度。}{0.96}{0.29\textheight}{R历史基准误差分布图.png}{附图 13 历史同长度窗口基准误差分布}

\appendixfigunit{0.60\textheight}{库存与供需缺口风险}{该图展示中长期外推中商业库存剩余率和剩余缺口风险的同步变化，说明悲观路径的尾部风险主要来自 SPR 等效预算耗尽、库存缓冲下降和剩余缺口同时存在。}{1.00}{0.51\textheight}{库存与供需缺口风险.png}{附图 14 库存与供需缺口风险}

上述补充图共同说明，长期外推中的随机扰动、风险变量与库存缺口并非孤立设定：历史窗口提供扰动尺度，GPR 与 OVX 提供外生风险约束，库存和剩余缺口则刻画尾部二次跳涨的物理来源。因此正文把长期结果写成条件路径、分位区间和尾部概率，而不是只给一条确定性均值预测线。

\clearpage
